\scanpage{71}
% archive sheet number: 72
5.\unsure{)} Puissances extérieures d'op. de Fredholm, inégalités diverses.

\uline{Prop. 1}

\[
\lVert u_1\wedge\ldots\wedge u_n\rVert_1\leq\frac1{n!}\lVert u_1\rVert_1\ldots\lVert u_n\rVert_1\tag*{\textcircled{1}}
\]

\uline{Corollaire 1}

\[
\left\{\begin{aligned}
\alpha_n(u_1,\ldots,u_n)&=\operatorname{Tr}(u_1\wedge\ldots\wedge u_n)\\
\alpha_n(u)&=\operatorname{Tr}\overset{n}{\Lambda}(u)
\end{aligned}\right.\tag*{$\substack{\text{\textcircled{2}}\\\text{\textcircled{3}}}$}
\]

\uline{Corollaire 2}

\[
|\alpha_n(u_1\deleted{\unsure{,\ldots|)}},\ldots,u_n)|\leq\frac1{n!}\lVert u_1\rVert_1\ldots\lVert u_n\rVert_1\quad\text{en particulier }\ |\alpha_n(u)|\leq\frac1{n!}\lVert u\rVert_1^n\tag*{(4)}
\]

\[
\deleted{(5)\quad|\alpha_n(u)|\leq\frac1{n!}\lVert u\rVert^{n}}
\]

\deleted{\uline{Corollaire 3}}

\[
\begin{aligned}
\deleted{(6)\quad\lVert\overset{n}{\Lambda}u\rVert_1}&\deleted{=\alpha_n(|u|)}\\
\deleted{(7)\quad\lVert\overset{n}{\Lambda}u\rVert_1}&\deleted{=\sqrt{\alpha_n(u^*u)}}
\end{aligned}
\]

\[
\deleted{\bigl(|\alpha_n(u)|\leq\lVert\overset{n}{\Lambda}u\rVert_1=\alpha_n(|u|)\leq\frac1{n!}\lVert u\rVert_1^n\bigr)}
\]

\deleted{\uline{Corollaire 4}}

\[
\deleted{(8)\quad\alpha_n(u)\leq\alpha_n(|u|)}
\]

\uline{Corollaire 3}

\[
|\alpha_n(u)|\leq\deleted{\unsure{\alpha_n|\overset{n}{\Lambda}}u}\alpha_n(|u|)=\lVert\overset{n}{\Lambda}u\rVert_1\leq\frac1{n!}\lVert u\rVert_1^n\tag*{\textcircled{5}}
\]

Car a $|\alpha_n(u)|\leq\lVert\overset{n}{\Lambda}u\rVert_1$ en vertu de (3), et $\alpha_n(|u|)=\operatorname{Tr}(\overset{n}{\Lambda}(|u|))=\operatorname{Tr}|\overset{n}{\Lambda}u|=\lVert\overset{n}{\Lambda}u\rVert_1$ (deux références) enfin $\lVert\overset{n}{\Lambda}u\rVert_1\leq\frac1{n!}\lVert u\rVert^n$ (formule (1))

\margin{?}

\deleted{\uline{Corollaire 4}}

\[
\deleted{(6)\quad\left\{\begin{aligned}
\lVert\overset{n}{\Lambda}u\rVert_1&=\alpha_n(|u|),\\
\lVert\overset{n}{\Lambda}u\rVert_{2}&=\sqrt{\alpha_n(u^*u)}
\end{aligned}\right.}
\]

\uline{Corollaire 4}

\[
|\det(1+u)|\leq\det(1+|u|)\tag*{\textcircled{6}}
\]

\uline{Théorème 1} Si $u\in L^1(E)$, alors \deleted{$\det(1+zu)$} $(\lambda_i(u))\in l^1$, et $\det(1+zu)$ est de \uline{genre} $0$ :

\[
\det(1+zu)=\prod_i(1+z\lambda_i(u))\tag*{\textcircled{7}}
\]

Équivaut à

\[
\begin{aligned}
\alpha_n(u)&=\sum_{i_1<\ldots<i_n}\lambda_{i_1}(u)\ldots\lambda_{\unsure{n i_n|i_n}}(u)\\
\operatorname{Tr}u&=\sum_i\lambda_i(u)
\end{aligned}\tag*{$\substack{\text{\textcircled{8}}\\\text{\textcircled{9}}}$}
\]

Implique

\[
\alpha_n(|u|)=\lVert\overset{n}{\Lambda}u\rVert_1=\sum_{i_1<\unsure{\ldots}<i_{\unsure{2|n}}}\rho_{i_1}(u)\ldots\rho_{i_n}(u)\tag*{\textcircled{10}}
\]

\margin{Ce qui revient \deleted{\unsure{au|un}} dont résulte de nouveau (8) en l'appliquant à $\overset{n}{\Lambda}u$.}

Pour la preuve, on \deleted{est} ramènera au cas où $u$ est hermitien positif, où c'est facile,

\uline{Prop 2}

\[
\deleted{\unsure{\lVert u_1\wedge v_1\wedge u\rVert} \lVert(\overset{p_1}{\Lambda}u_1)\wedge\ldots\wedge(\overset{p_k}{\Lambda}u_k)\rVert_1\leq\lVert\overset{p_1}{\Lambda}u_1\rVert_1\ldots\lVert\overset{p_k}{\Lambda}u_k\rVert}
\]

\[
\lVert(\overset{p}{\Lambda}u)\wedge(\overset{q}{\Lambda}v)\rVert_1\leq\frac{p!q!}{(p+q)!}\lVert\overset{\unsure{p|0}}{\Lambda}u\rVert_1\lVert\overset{q}{\Lambda}v\rVert_1=\frac{p!q!}{(p+q)!}\alpha_p(|u|)\alpha_q(|v|)\tag*{\textcircled{11}}
\]

(le même vrai pour plus de deux facteurs, l'écriture $(\overset{p}{\Lambda}u)\wedge(\overset{q}{\Lambda}v)$ est ici \uline{mnémotechnique}).
